\section{Giới thiệu đề tài}
\label{sec:introduction}

Ngày nay, gần như mọi ứng dụng mà chúng ta sử dụng hằng ngày, từ ứng dụng ngân hàng, trang thương mại điện tử cho đến mạng xã hội, đều được xây dựng xoay quanh dữ liệu. Mỗi lần người dùng đăng nhập, tìm kiếm một sản phẩm hay đọc một bài viết, phía sau đều có hàng loạt thao tác lưu trữ và truy xuất dữ liệu diễn ra. Để những thao tác đó được thực hiện nhanh, chính xác và an toàn, các ứng dụng không tự quản lý dữ liệu trên các tệp tin rời rạc mà giao phó công việc này cho một \textbf{hệ quản trị cơ sở dữ liệu} (\emph{Database Management System} -- DBMS). Có thể hiểu DBMS là lớp phần mềm đứng giữa ứng dụng và dữ liệu: ứng dụng chỉ cần mô tả dữ liệu nào cần lưu hoặc cần lấy, còn việc tổ chức dữ liệu trên đĩa hay trên bộ nhớ, tìm ra cách truy xuất hiệu quả, giữ cho dữ liệu đúng khi nhiều người cùng thao tác và bảo vệ dữ liệu khi hệ thống gặp sự cố đều do DBMS đảm nhận. Nhờ vậy, người phát triển có thể tập trung vào nghiệp vụ của ứng dụng thay vì phải tự giải quyết những bài toán phức tạp về lưu trữ.

Tuy nhiên, không có một DBMS nào phù hợp cho mọi bài toán. Mỗi hệ quản trị cân bằng giữa các yêu cầu như tính toàn vẹn, khả năng truy vấn linh hoạt, tốc độ và độ bền của dữ liệu theo một cách riêng, và sự lựa chọn đó thể hiện rõ trong thiết kế bên trong của nó. Hiện nay có hai hướng tiếp cận tiêu biểu là các hệ quản trị cơ sở dữ liệu quan hệ (SQL) và các hệ NoSQL. Trong bài tập lớn này, nhóm chọn một đại diện cho mỗi hướng là \textbf{Microsoft SQL Server} và \textbf{Redis}. SQL Server là hệ quản trị cơ sở dữ liệu quan hệ được sử dụng rộng rãi trong doanh nghiệp: dữ liệu được tổ chức thành các bảng có lược đồ chặt chẽ, lưu trên đĩa và được truy vấn bằng ngôn ngữ T-SQL; bên cạnh đó, hệ quản trị này còn cung cấp hệ thống chỉ mục đa dạng, bộ tối ưu truy vấn, giao tác bảo đảm các tính chất ACID cùng cơ chế sao lưu và phục hồi đầy đủ. Ngược lại, Redis là một hệ lưu trữ khoá -- giá trị đặt toàn bộ dữ liệu trong bộ nhớ chính. Redis không có khái niệm lược đồ mà cung cấp các cấu trúc dữ liệu như chuỗi, bảng băm, danh sách, tập hợp và tập hợp có thứ tự, nhờ đó mỗi thao tác chỉ mất thời gian rất ngắn. Chính vì vậy, Redis thường được dùng làm bộ nhớ đệm, bộ đếm hoặc nơi lưu phiên đăng nhập. Có thể thấy hai hệ quản trị này đi theo hai triết lý gần như đối lập nhau: SQL Server ưu tiên sự toàn vẹn và độ bền của dữ liệu, trong khi Redis ưu tiên tốc độ. Cũng vì sự khác biệt đó mà trong thực tế, hai loại hệ quản trị này thường không thay thế nhau mà được kết hợp trong cùng một hệ thống, trong đó cơ sở dữ liệu quan hệ giữ vai trò nguồn dữ liệu gốc còn Redis đóng vai trò lớp phục vụ nhanh đứng phía trước.

Xuất phát từ những nhận định trên, bài tập lớn hướng tới hai mục tiêu gắn liền với nhau. Mục tiêu thứ nhất là tìm hiểu và giải thích các cơ chế bên trong của SQL Server và Redis, từ đó so sánh hai hệ quản trị một cách có cơ sở. Cụ thể, nhóm nghiên cứu năm chủ đề kỹ thuật, tương ứng với năm thành viên, gồm: quản lý và lưu trữ dữ liệu, cấu trúc chỉ mục, xử lý truy vấn, xử lý giao tác, sao lưu và phục hồi dữ liệu. Với mỗi chủ đề, nhóm trình bày cơ sở lý thuyết, phân tích cách mỗi hệ quản trị hiện thực vấn đề đó, rồi thiết kế các kịch bản thực nghiệm tương đương trên cả hai hệ để kiểm chứng và đối chiếu kết quả. Mục tiêu thứ hai là vận dụng những hiểu biết này vào việc xây dựng một ứng dụng thực tế. Ứng dụng được nhóm lựa chọn là \textbf{Novelle}, một trang web đọc và sáng tác truyện trực tuyến dành cho bốn nhóm người dùng: khách, người đọc, tác giả và quản trị viên. Đây là một bối cảnh phù hợp để kết hợp hai hệ quản trị, bởi vì dữ liệu của một nền tảng truyện vừa có nhiều quan hệ chặt chẽ cần được bảo đảm toàn vẹn, vừa có lượng truy cập đọc rất lớn cần được phục vụ nhanh. Trong Novelle, SQL Server lưu toàn bộ dữ liệu nghiệp vụ như tài khoản, truyện, chương và đánh giá, còn Redis đảm nhận bộ nhớ đệm, bộ đếm lượt xem, bản nháp, phiên đăng nhập và giới hạn tần suất. Cách phân chia vai trò này không được chọn theo cảm tính mà được giải thích dựa trên chính kết quả của các chủ đề nghiên cứu.

Để trình bày các nội dung trên một cách có hệ thống, phần còn lại của báo cáo được tổ chức như sau. Chương~\ref{sec:dbms-overview} giới thiệu tổng quan về kiến trúc và đặc điểm nổi bật của SQL Server và Redis, qua đó đưa ra một so sánh ban đầu giữa hai hệ quản trị. Tiếp theo, Chương~\ref{sec:app-db-design} mô tả ứng dụng Novelle cùng các yêu cầu chức năng, sau đó trình bày quá trình mô hình hoá dữ liệu, thiết kế lược đồ quan hệ trên SQL Server và cách tổ chức dữ liệu trên Redis. Chương~\ref{sec:dbms-topics} là phần trọng tâm của báo cáo, lần lượt trình bày năm chủ đề nghiên cứu, mỗi chủ đề gồm cơ sở lý thuyết, cơ chế trên từng hệ quản trị, thực nghiệm và đánh giá, và khép lại bằng phần kết luận chung cho cả hai hệ quản trị. Dựa trên những kết quả đó, Chương~\ref{sec:architecture-implementation} trình bày các công nghệ được sử dụng, kiến trúc hệ thống, cách tích hợp hai hệ quản trị và quá trình hiện thực các chức năng chính của ứng dụng. Cuối cùng, Chương~\ref{sec:conclusion} tổng kết những kết quả đã đạt được, đánh giá mức độ hoàn thành và chỉ ra những hạn chế còn tồn tại của đề tài. Danh mục tài liệu tham khảo và phụ lục chứa biên bản họp nhóm được đặt ở cuối báo cáo.
